\documentclass[11pt,a4paper]{article}

\usepackage[T1]{fontenc}
\usepackage[utf8]{inputenc}
\usepackage[english]{babel}
\usepackage[a4paper,top=2.35cm,bottom=2.45cm,left=2.55cm,right=2.55cm,headsep=12pt]{geometry}
\usepackage{amsmath,amssymb,mathtools}
\usepackage{enumitem}
\usepackage{eurosym}
\usepackage{graphicx}
\usepackage{xcolor}
\usepackage{microtype}
\usepackage{avant}
\usepackage{mathptmx}
\usepackage{titlesec}
\usepackage{fancyhdr}
\usepackage{mdframed}
\usepackage{tikz}
\usepackage{needspace}
\usepackage{hyperref}
\usetikzlibrary{arrows.meta,positioning}

\definecolor{PPEBurgundy}{RGB}{112,28,47}
\definecolor{PPENavarraRed}{RGB}{211,32,46}
\definecolor{PPEGold}{RGB}{185,143,57}
\definecolor{PPEAccentLight}{RGB}{248,235,238}
\definecolor{PPEInk}{RGB}{37,31,33}
\definecolor{PPEMuted}{RGB}{101,83,88}

\hypersetup{
  colorlinks=true,
  linkcolor=PPEBurgundy,
  urlcolor=PPEBurgundy
}

\setlength{\parindent}{0pt}
\setlength{\parskip}{0.45em}
\setlist[itemize]{topsep=0.3em,itemsep=0.2em,leftmargin=1.5em}
\setlist[enumerate]{topsep=0.35em,itemsep=0.35em,leftmargin=1.9em}
\raggedbottom
\allowdisplaybreaks[2]

\titleformat{\section}
  {\sffamily\bfseries\large\color{PPEInk}}
  {\color{PPEBurgundy}\thesection}{0.8em}{}
\titlespacing*{\section}{0pt}{1.45em}{0.55em}

\newmdenv[
  skipabove=0.7em,
  skipbelow=0.9em,
  leftline=true,rightline=false,topline=false,bottomline=false,
  linewidth=3pt,
  linecolor=PPEBurgundy,
  backgroundcolor=PPEAccentLight,
  innerleftmargin=0.8em,innerrightmargin=0.8em,
  innertopmargin=0.55em,innerbottommargin=0.55em
]{solutionbox}

\newmdenv[
  skipabove=0.65em,
  skipbelow=0.75em,
  leftline=true,rightline=false,topline=false,bottomline=false,
  linewidth=2pt,
  linecolor=PPEGold,
  backgroundcolor=white,
  innerleftmargin=0.75em,innerrightmargin=0.75em,
  innertopmargin=0.45em,innerbottommargin=0.45em
]{tipbox}

\pagestyle{fancy}
\fancyhf{}
\fancyhead[LE,LO]{\sffamily\small Microeconomics \textbar{} Exercise Set I \textbar{} Worked solutions}
\fancyhead[RE,RO]{\sffamily\small\color{PPEBurgundy}\thepage}
\fancyfoot[C]{\sffamily\scriptsize\color{PPEMuted}Companion to Exercise Set I, Version B}
\renewcommand{\headrulewidth}{0.45pt}
\renewcommand{\headrule}{\hbox to\headwidth{\color{PPEGold}\leaders\hrule height \headrulewidth\hfill}}
\setlength{\headheight}{15pt}

\newcommand{\MRS}{\operatorname{MRS}}
\newcommand{\MU}{\operatorname{MU}}

\begin{document}

\begin{center}
  \includegraphics[width=0.28\textwidth]{assets/branding/unav-logo-black.png}\par
  \vspace{0.55em}
  {\sffamily\color{PPENavarraRed}\large\textsc{Politics, Philosophy and Economics}\par}
  \vspace{0.3em}
  {\color{PPEGold}\rule{0.72\textwidth}{1.3pt}\par}
  \vspace{0.55em}
  {\sffamily\bfseries\color{PPEBurgundy}\LARGE Microeconomics: Exercise Set I\par}
  \vspace{0.16em}
  {\sffamily\large Worked solutions --- Version B --- Lectures 1--4\par}
\end{center}

\begin{solutionbox}
\textbf{How to use this key.} Try each question before reading its solution.
The worked answers show the set-up, calculation, conclusion, and the economic
interpretation. Equivalent algebra and clearly explained diagrams are also
valid answers.
\end{solutionbox}

\section{Mathematical tools}

\subsection*{1.1 Equations}
\begin{align*}
\text{(1)}\quad 420-7x&=0
&&\Longrightarrow& 7x&=420
&&\Longrightarrow& \boxed{x=60}.\\[0.3em]
\text{(2)}\quad \frac{144}{x}&=9
&&\Longrightarrow& 144&=9x
&&\Longrightarrow& \boxed{x=16}.\\[0.3em]
\text{(3)}\quad (x-7)^2&=0
&&\Longrightarrow& \boxed{x=7}.\\[0.3em]
\text{(4)}\quad (x-3)^2&=49
&&\Longrightarrow& x-3&=\pm7
&&\Longrightarrow& \boxed{x=10\text{ or }x=-4}.\\[0.3em]
\text{(5)}\quad 8x-72+25x&=324
&&\Longrightarrow& 33x&=396
&&\Longrightarrow& \boxed{x=12}.
\end{align*}

\subsection*{1.2 Derivatives}
\textbf{(6)} A derivative is the instantaneous rate at which a function
changes with its argument:
\[
  f'(x)=\lim_{h\to0}\frac{f(x+h)-f(x)}{h},
\]
when this limit exists. Geometrically it is the slope of the tangent line.
For example, if \(C(q)\) is total cost, then \(C'(q)\) is marginal cost:
the approximate additional cost of producing one more unit near output \(q\).

\begin{align*}
\text{(7)}\quad f(x)=x^3
&&\Longrightarrow& f'(x)&=\boxed{3x^2}.\\[0.3em]
\text{(8)}\quad f(x)=\frac25x^4+5x^2-8x
&&\Longrightarrow& f'(x)&=\boxed{\frac85x^3+10x-8}.\\[0.3em]
\text{(9)}\quad f(x)=\frac25x^4+5x^2-8x+430
&&\Longrightarrow& f'(x)&=\boxed{\frac85x^3+10x-8}.
\end{align*}
The constant \(430\) disappears because the derivative of a constant is zero.
\begin{align*}
\text{(10)}\quad f(x)&=3x^{18}-4x^{12}+7x^6-9x+100\\
&&\Longrightarrow\quad f'(x)&=\boxed{54x^{17}-48x^{11}+42x^5-9}.\\[0.3em]
\text{(11)}\quad f(x,y)&=3x^2+4xy+y^2-7
&&\Longrightarrow\quad \frac{\partial f}{\partial x}
  &=\boxed{6x+4y}.\\[0.3em]
\text{(12)}\quad f(x,y)&=x^3y-2xy+5x
&&\Longrightarrow\quad \frac{\partial f}{\partial x}
  &=\boxed{3x^2y-2y+5}.\\[0.3em]
\text{(13)}\quad f(x,y)&=x^{2/5}y^{3/5}
&&\Longrightarrow\quad \frac{\partial f}{\partial x}
  &=\boxed{\frac25x^{-3/5}y^{3/5}}.
\end{align*}
In (11)--(13), \(y\) is held fixed while differentiating with respect to
\(x\). In (13), the displayed formula applies for \(x>0\); at \(x=0\) it
is not finite when \(y>0\).

\section{Preference language and monotonicity}

\textbf{2.a Preference notation.} The statement \(A\sim B\) means that the
consumer is indifferent between bundles \(A\) and \(B\): each is at least as
good as the other. The statement \(A\succeq B\succ C\) means that \(A\) is
weakly preferred to \(B\), and \(B\) is strictly preferred to \(C\). With
transitive preferences, it also follows that \(A\succ C\).

\textbf{2.b Comparing the bundles.} The diagram gives
\[
  A=(1,2),\qquad B=(5,2),\qquad C=(1,6),\qquad D=(5,6).
\]
For ordinary goods, a bundle with at least as much of both goods and more of
one is strictly preferred. Therefore,
\[
  \boxed{B\succ A,\qquad C\succ A,\qquad D\succ B,\qquad D\succ C.}
\]
The comparison between \(B\) and \(C\) cannot be determined from monotonicity
alone: \(B\) has more \(x\), while \(C\) has more \(y\). The consumer's
trade-off between the goods is needed; \(B\succ C\), \(C\succ B\), and
\(B\sim C\) are all possible.

\begin{tipbox}
\textbf{Useful distinction.} Monotonicity compares bundles only when one
bundle has weakly more of every ordinary good. It does not tell us how a
consumer trades one good for another.
\end{tipbox}

\section{Utility, budget constraints, and consumer choice}

\needspace{10\baselineskip}
\textbf{3.a Elena's choice.}

\textit{Set-up.} The budget set and its boundary are
\[
  B=\{(x,y)\in\mathbb{R}_+^2:6x+3y\leq72\},
  \qquad
  6x+3y=72,
\]
or, solving for \(y\),
\[
  y=24-2x.
\]
The intercepts are \((12,0)\) and \((0,24)\); the slope is
\(-p_x/p_y=-2\). The indifference curves are
\[
  xy=\bar u
  \qquad\Longleftrightarrow\qquad
  y=\frac{\bar u}{x}.
\]
Thus the requested curves are \(y=36/x\), \(y=54/x\), and \(y=72/x\).

\textit{Optimisation by substitution.} On the budget line,
\[
  u(x) = x(24-2x)=24x-2x^2.
\]
Hence
\[
  u'(x)=24-4x=0
  \quad\Longrightarrow\quad x^*=6,
  \qquad
  u''(x)=-4<0.
\]
The negative second derivative shows that this stationary point is a maximum.
The budget line gives
\[
  y^*=24-2(6)=12,
  \qquad
  u(x^*,y^*)=6\cdot12=72.
\]
The marginal utilities are \(\MU_x=y\) and \(\MU_y=x\), so
\[
  \MRS_{x,y}=\frac{\MU_x}{\MU_y}=\frac{y}{x}.
\]
At \((6,12)\), \(\MRS_{x,y}=12/6=2=p_x/p_y=6/3\), which verifies
tangency.

\begin{center}
\begin{tikzpicture}[x=0.43cm,y=0.23cm,>=Latex]
  \draw[->,thick] (0,0) -- (13.2,0) node[right] {\(x\): tickets};
  \draw[->,thick] (0,0) -- (0,25.8) node[above] {\(y\): lunches};
  \draw[PPEBurgundy,very thick] (0,24) -- (12,0)
    node[pos=0.22,above right,black] {budget line};
  \draw[PPEGold,thick,domain=1.6:11.5,samples=80,smooth]
    plot (\x,{72/\x});
  \node[PPEGold,above] at (8.2,8.8) {\(u=72\)};
  \filldraw[PPENavarraRed] (6,12) circle (2.4pt)
    node[above right,black] {\((6,12)\)};
  \draw[densely dashed,PPEMuted] (6,0) -- (6,12) -- (0,12);
  \node[below] at (6,0) {\(6\)};
  \node[left] at (0,12) {\(12\)};
\end{tikzpicture}
\end{center}

\begin{solutionbox}
\textbf{Result for Exercise 3.} Elena chooses
\[
  \boxed{(x^*,y^*)=(6,12),\qquad u^*=72.}
\]
\end{solutionbox}

\needspace{9\baselineskip}
\textbf{4.a Choice with \(u(x,y)=y(x+3)\).}

\textit{Set-up.} The budget line is
\[
  6x+7y=66
  \qquad\Longleftrightarrow\qquad
  y=\frac{66-6x}{7}.
\]
Its intercepts are \((11,0)\) and \((0,66/7)\). The requested indifference
curves are \(y=21/(x+3)\) and \(y=42/(x+3)\).

The marginal utilities are
\[
  \MU_x=y,
  \qquad
  \MU_y=x+3,
  \qquad
  \MRS_{x,y}=\frac{y}{x+3}.
\]
At an interior optimum,
\[
  \frac{y}{x+3}=\frac{p_x}{p_y}=\frac67,
  \qquad\text{so}\qquad
  y=\frac67(x+3).
\]
Substitution into the budget line gives
\[
  6x+7\left[\frac67(x+3)\right]=66
  \Longrightarrow
  12x+18=66
  \Longrightarrow
  x^*=4.
\]
Then \(y^*=6\), and utility is \(u^*=6(4+3)=42\).

\textit{Verification by substitution.}
\[
  u(x)=\frac{66-6x}{7}(x+3)
  =\frac{-6x^2+48x+198}{7}.
\]
Therefore,
\[
  u'(x)=\frac{-12x+48}{7}=0
  \Longrightarrow x=4,
  \qquad
  u''(x)=-\frac{12}{7}<0.
\]
This confirms a maximum.

\begin{solutionbox}
\textbf{Result for Exercise 4.}
\[
  \boxed{(x^*,y^*)=(4,6),\qquad u^*=42,\qquad
  \MRS_{x,y}=\frac67.}
\]
\end{solutionbox}

\section{Rational preferences}

\textbf{5.a The axioms.} Let \(A\), \(B\), and \(C\) be feasible bundles.

\begin{description}[style=nextline,font=\sffamily\bfseries,leftmargin=2.65cm]
  \item[Completeness]
  The consumer can compare every pair: \(A\succeq B\) or \(B\succeq A\)
  (or both). For example, the student can say whether a coffee-and-sandwich
  bundle is better than, worse than, or as good as a tea-and-cake bundle.

  \item[Transitivity]
  If \(A\succeq B\) and \(B\succeq C\), then \(A\succeq C\). It rules out
  cycles: ranking tacos at least as highly as pizza and pizza at least as
  highly as salad requires ranking tacos at least as highly as salad.

  \item[Reflexivity]
  Every bundle is at least as good as itself: \(A\succeq A\). The lunch
  bundle a student has today is, trivially, at least as good as that same
  lunch bundle.

  \item[Monotonicity]
  If \(A\) has at least as much of every ordinary good as \(B\), and more of
  at least one, then \(A\succ B\). Holding everything else fixed, an extra
  museum ticket or an extra lunch makes the bundle strictly better.

  \item[Continuity]
  Preference rankings do not reverse after an arbitrarily small change in a
  bundle. Formally, if \(A_n\succeq B_n\) for every \(n\), \(A_n\to A\), and
  \(B_n\to B\), then \(A\succeq B\). A tiny change in the size of a meal
  should not create an abrupt, unexplained reversal of a ranking.

  \item[Convexity]
  If \(A\succeq C\) and \(B\succeq C\), then every mixture
  \(tA+(1-t)B\), with \(0\leq t\leq1\), is also at least as good as \(C\).
  It captures a taste for variety: if two weekly food bundles are both
  acceptable alternatives to a third, a balanced combination is also
  acceptable.
\end{description}

\textbf{5.b The food ranking.} The ranking violates \textbf{transitivity}.
From tacos \(\succeq\) pizza and pizza \(\succeq\) salad, transitivity
requires tacos \(\succeq\) salad. But salad \(\succ\) tacos means that tacos
are not at least as good as salad. The three statements cannot all describe
one transitive ordering.

\section{Budget sets, relative prices, and opportunity cost}

\textbf{6.a Campus lunches and drinks.}

\begin{enumerate}[label=(\roman*)]
  \item The budget set and budget line are
  \[
    B=\{(x,y)\in\mathbb{R}_+^2:6x+3y\leq30\},
    \qquad
    y=10-2x.
  \]
  The intercepts are \(\boxed{(5,0)}\) and \(\boxed{(0,10)}\).

  \item One lunch costs \euro6 and one drink costs \euro3. One additional
  lunch therefore costs
  \[
    \frac{p_x}{p_y}=\frac63=\boxed{2\text{ drinks}}.
  \]
  This is also the absolute value of the budget-line slope.

  \item With \(I=42\), the new line is
  \[
    6x+3y=42
    \qquad\Longleftrightarrow\qquad
    y=14-2x.
  \]
  Both intercepts increase to \((7,0)\) and \((0,14)\). With prices
  unchanged, the budget line shifts outward in parallel.

  \item When \(p_y=6\), the new line is
  \[
    6x+6y=42
    \qquad\Longleftrightarrow\qquad
    y=7-x.
  \]
  The \(x\)-intercept remains \(7\), whereas the \(y\)-intercept falls from
  \(14\) to \(7\). The line pivots inward around the \(x\)-intercept and
  becomes flatter, changing from slope \(-2\) to slope \(-1\).
\end{enumerate}

\section{Consumer maximisation and a price change}

\needspace{9\baselineskip}
\textbf{7.a Initial consumer problem.} The budget line is
\[
  2x+y=12
  \qquad\Longleftrightarrow\qquad
  y=12-2x.
\]
The marginal utilities are
\[
  \MU_x=2+y,
  \qquad
  \MU_y=1+x,
  \qquad
  \MRS_{x,y}=\frac{2+y}{1+x}.
\]
The tangency condition is
\[
  \frac{2+y}{1+x}=\frac{p_x}{p_y}=2
  \Longrightarrow y=2x.
\]
Together with \(2x+y=12\), this yields
\[
  4x=12
  \Longrightarrow
  x^*=3,
  \qquad
  y^*=6.
\]
For a direct check, substitute the budget line:
\[
  u(x)=2x+(12-2x)+x(12-2x)=12+12x-2x^2.
\]
Thus
\[
  u'(x)=12-4x=0
  \Longrightarrow x=3,
  \qquad
  u''(x)=-4<0.
\]
The stationary point is therefore a maximum; indeed,
\(\MRS_{x,y}=8/4=2\) at \((3,6)\).

\begin{solutionbox}
\textbf{Result for Exercise 7.}
\[
  \boxed{(x^*,y^*)=(3,6),\qquad u^*=30.}
\]
\end{solutionbox}

\needspace{9\baselineskip}
\textbf{8.a Higher price of \(x\).} The new budget line is
\[
  3x+y=12
  \qquad\Longleftrightarrow\qquad
  y=12-3x.
\]
The tangency condition now becomes
\[
  \frac{2+y}{1+x}=3
  \Longrightarrow y=1+3x.
\]
Combining this condition with the budget line gives
\[
  3x+(1+3x)=12
  \Longrightarrow
  x^*=\frac{11}{6},
  \qquad
  y^*=\frac{13}{2}.
\]
Equivalently,
\[
  u(x)=2x+(12-3x)+x(12-3x)=12+11x-3x^2,
\]
so \(u'(x)=11-6x=0\) and \(u''(x)=-6<0\).

The line rotates inward around the \(y\)-intercept: the \(y\)-intercept
remains \(12\), but the \(x\)-intercept falls from \(6\) to \(4\), and the
slope changes from \(-2\) to \(-3\). The quantity of \(x\) falls from \(3\)
to \(11/6\).

This is a \textbf{movement along} the demand curve for \(x\), because
\(p_x\) itself changed while income, preferences, and \(p_y\) were fixed.
Changes in income or in the price of the related good \(y\), for example,
would shift the demand curve for \(x\) instead.

\begin{tipbox}
\textbf{Common terminology error.} A change in the good's own price moves
the consumer along its demand curve. A shift requires another determinant of
demand to change.
\end{tipbox}

\begin{solutionbox}
\textbf{Result for Exercise 8.}
\[
  \boxed{(x^*,y^*)=\left(\frac{11}{6},\frac{13}{2}\right).}
\]
\end{solutionbox}

\section{MRS, convexity, and demand classifications}

\textbf{9.a MRS and convexity.} For \(u(x,y)=y(x+3)\),
\[
  \MU_x=y,
  \qquad
  \MU_y=x+3,
  \qquad
  \MRS_{x,y}=\frac{\MU_x}{\MU_y}=\frac{y}{x+3}.
\]
The MRS is the amount of \(y\) the consumer is willing to give up for one
additional unit of \(x\), while remaining on the same indifference curve.
The slope of an indifference curve is \(-\MRS_{x,y}\).

Along \(u(x,y)=\bar u>0\),
\[
  y=\frac{\bar u}{x+3}.
\]
Substitution into the MRS gives
\[
  \MRS_{x,y}
  =\frac{\bar u/(x+3)}{x+3}
  =\boxed{\frac{\bar u}{(x+3)^2}}.
\]
Moreover,
\[
  \frac{d\MRS_{x,y}}{dx}
  =-\frac{2\bar u}{(x+3)^3}<0.
\]
Thus the MRS diminishes as \(x\) rises: the indifference curve becomes
flatter. This is the smooth graphical expression of convex preferences and
a taste for balanced bundles.

\textbf{9.b Normal and inferior goods.} Holding prices fixed, a good is
\textbf{normal} if \(\partial q/\partial I>0\), and \textbf{inferior} if
\(\partial q/\partial I<0\). Here,
\[
  \frac{\partial q_{\mathrm{bus}}}{\partial I}=-0.02<0,
  \qquad
  \frac{\partial q_{\mathrm{coffee}}}{\partial I}=0.015>0.
\]
Therefore, bus trips are \(\boxed{\text{inferior}}\), and coffee is
\(\boxed{\text{normal}}\) (also called superior in some introductory
terminology). Both quantities remain non-negative over the stated income
range.

\textbf{9.c Why a Giffen good is inferior.} When the price of \(x\) falls,
the substitution effect always encourages consumption of \(x\), because it
has become relatively cheaper. A Giffen good is one for which total demand
nevertheless falls when its price falls. Therefore the negative income
effect of the price fall must be larger than the positive substitution
effect. An income effect that reduces consumption as real income rises is
the defining feature of an inferior good. Hence every Giffen good is
inferior, although most inferior goods are not Giffen goods.

\section{Aggregating individual demands}

\textbf{10.a Market demand.} Group \(A\) buys while \(p<8\), and group \(B\)
buys while \(p<10\). Add their quantities at each common price:
\[
  Q(p)=
  \begin{cases}
    13-\dfrac32p, & 0\leq p<8,\\[0.45em]
    5-\dfrac12p, & 8\leq p<10,\\[0.45em]
    0, & p\geq10.
  \end{cases}
\]

\textbf{10.b Choke prices and kink.} Group \(A\)'s choke price is
\(\boxed{\euro8}\), and group \(B\)'s is \(\boxed{\euro10}\). At
\(p=8\), group \(A\) stops buying but group \(B\) still demands one bottle.
Thus \(p=8\) is the interior kink. The aggregate choke price is
\(\boxed{\euro10}\), where market demand reaches zero.

\textbf{10.c Sketch and interpretation.} With quantity on the horizontal
axis and price on the vertical axis, the market curve passes through
\[
  (Q,p)=(13,0),\qquad (1,8),\qquad (0,10).
\]
\begin{center}
\begin{tikzpicture}[x=0.48cm,y=0.48cm,>=Latex]
  \draw[->,thick] (0,0) -- (14.3,0) node[right] {quantity \(Q\)};
  \draw[->,thick] (0,0) -- (0,11.5) node[above] {price \(p\)};
  \draw[PPEBurgundy,very thick] (13,0) -- (1,8) -- (0,10);
  \filldraw[PPENavarraRed] (13,0) circle (2.2pt)
    node[below right,black] {\((13,0)\)};
  \filldraw[PPENavarraRed] (1,8) circle (2.2pt)
    node[above right,black] {\((1,8)\)};
  \filldraw[PPENavarraRed] (0,10) circle (2.2pt)
    node[above left,black] {\((0,10)\)};
  \draw[densely dashed,PPEMuted] (1,0) -- (1,8) -- (0,8);
  \node[below] at (1,0) {\(1\)};
  \node[left] at (0,8) {\(8\)};
  \node[left] at (0,10) {\(10\)};
\end{tikzpicture}
\end{center}

We add quantities horizontally because all consumers face the same market
price, and the total number of bottles demanded is the sum of their
quantities. Adding prices would not make economic sense: price is one common
unit price, whereas quantities are the objects that can be added across
buyers.

\begin{solutionbox}
\textbf{Result for Exercise 10.} The market demand curve has an interior
kink at \(p=8\), and its final choke price is \(p=10\).
\end{solutionbox}

\vfill
\begin{center}
  {\sffamily\small\color{PPEMuted}
  Worked solutions for Exercise Set I, Version B.}
\end{center}

\end{document}
